% ==========================================
%  content/05-declaration.tex  —  成果、致谢、原创性声明
%  可独立编译: xelatex 05-declaration.tex
% ==========================================
\documentclass[../main.tex]{subfiles}
\begin{document}

% ==========================================
%  攻读学位期间完成的论文
%  ★ 如果不需要此章节（如本科论文），删除下面直到“致谢”之前的所有行
% ==========================================
\cleardoublepage
\begin{center}
	{\heiti \zihao{3} 攻读\youreduction 士学位期间完成的论文}
\end{center}
\vspace{0.1cm}
\songti \zihao{4}
\phantomsection\addcontentsline{toc}{section}{\heiti \zihao{-4} 攻读\youreduction 士学位期间完成的论文}\tolerance=500

{
	\zihao{-4}
	\begin{enumerate}[label=\textbf{[\arabic*]}]
		% 按格式填写论文：
		% \item 作者1, 作者2, \& 作者3. 论文标题. 期刊名 卷号, 页码 (年份).
		\item First Author, Second Author, \& Third Author.
			Paper title here.
			\textit{Journal Name} \textbf{Volume}, Page (Year).
		\item ...
	\end{enumerate}
}

% ==========================================
%  致谢
% ==========================================
\cleardoublepage
\begin{center}
	{\heiti \zihao{3} 致\hspace{2em}谢}
\end{center}
\vspace{0.5cm}
\phantomsection\addcontentsline{toc}{section}{\heiti \zihao{-4} 致\hspace{2em}谢}\tolerance=500

\begin{spacing}{1.5}
	\songti \zihao{-4}

	在此处填写致谢内容。

	\rightline{\kaishu \zihao{-3} \yourname}
	\rightline{\yourtitledate}
\end{spacing}

% ==========================================
%  原创性声明与版权使用授权书
% ==========================================
\cleardoublepage
\newgeometry{left=3.4cm,right=3.4cm,top=4.5cm,bottom=2.5cm,headsep=25pt}
\thispagestyle{empty}

\begin{spacing}{1.5}

	% --- 原创性声明 ---
	\begin{center}
		\zihao{3} \heiti
		\universityname 学位论文原创性声明
	\end{center}
	\phantomsection\addcontentsline{toc}{section}{\heiti \zihao{-4} \universityname 学位论文原创性声明}\tolerance=500

	\begin{flushleft}
		\zihao{3} \songti
		\qquad
		本人郑重声明：所呈交的学位论文，是本人在导师的指导下，独立进行研究工作所取得的成果。
		除文中已经注明引用的内容外，本论文不含任何其他个人或集体已经发表或撰写过的作品成果。
		对本文的研究做出重要贡献的个人和集体，均已在文中以明确的方式标明。
		本人完全意识到本声明的法律结果由本人承担。
	\end{flushleft}

	\begin{center}
		\begin{multicols}{2}
			\balance
			\songti \zihao{3}
			\makebox[8em][s]{学位论文作者签名}:  \\[8em]
			\makebox[8em][s]{\hspace{\fill}年\hspace{\fill}月\hspace{\fill}日} \\[0em]
		\end{multicols}
	\end{center}
	\vskip 1cm

	% --- 版权使用授权书 ---
	\begin{center}
		\zihao{3} \heiti
		\universityname 学位论文版权使用授权书
	\end{center}
	\phantomsection\addcontentsline{toc}{section}{\heiti \zihao{-4} \universityname 学位论文版权使用授权书}\tolerance=500

	\begin{flushleft}
		\zihao{3} \songti
		\qquad
		本学位论文作者完全了解学校有关保留、使用学位论文的规定，
		研究生在校攻读学位期间论文工作的知识产权单位属\universityname。
		同意学校保留并向国家有关部门或机构送交论文的复印件和电子版，
		允许论文被查阅和借阅。本人授权\universityname 可以将本学位论文的全部或部分内容
		编入有关数据库进行检索，可以采用影印、缩印或扫描等复制手段保存和汇编本学位论文。
	\end{flushleft}
	\vskip 0.5em

	\begin{center}
		\begin{multicols}{2}
			\balance
			\songti \zihao{3}
			\makebox[5em][s]{作者签名：}  \\
			\makebox[5em][s]{导师签名：}  \\
			\makebox[5em][s]{日期：\hspace{\fill}} 年 \quad 月 \quad 日 \\
			\makebox[5em][s]{日期：\hspace{\fill}} 年 \quad 月 \quad 日 \\
		\end{multicols}
	\end{center}

\end{spacing}

\restoregeometry

\end{document}
